% Part: methods
% Chapter: induction
% Section: inductive-definitions

\documentclass[../../../include/open-logic-section]{subfiles}

\begin{document}

\olfileid[pt]{mth}{ind}{idf}

\olsection{Definições Indutivas}

Em lógica, muito frequentemente definimos tipos de objetos \emph{indutivamente},
isto é, especificando regras para o que conta como um objeto do tipo a
ser definido, que explicam como obter novos objetos desse tipo a partir de objetos
antigos desse tipo.  Por exemplo, frequentemente definimos tipos especiais de
sequências de símbolos, como os termos e !!{formula}s de uma
linguagem, por indução.  Para um exemplo simples, considere cadeias
compostas das letras $\mathrm{a}$, $\mathrm{b}$, $\mathrm{c}$,
$\mathrm{d}$, do símbolo~$\circ$, e dos colchetes $[$ e $]$, como
``$[[\mathrm{c} \circ \mathrm{d}][$'', ``$[\mathrm{a}[]\circ]$'',
``$\mathrm{a}$'' ou ``$[[\mathrm{a} \circ \mathrm{b}]\circ
\mathrm{d}]$''.  Você provavelmente sente que há algo
``errado'' com as duas primeiras cadeias: os colchetes não
``se equilibram'' de modo algum na primeira, e você pode sentir que o
``$\circ$'' deveria ``conectar'' expressões que elas próprias façam
sentido. A terceira e a quarta cadeia parecem melhores: para todo ``$[$''
há um ``$]$'' de fechamento (se houver algum), e para qualquer
$\circ$ podemos encontrar expressões ``bonitas'' de cada lado,
cercadas por um par de parênteses.

Gostaríamos de especificar precisamente o que conta como um ``termo bonito.''
Em primeiro lugar, toda letra isolada é bonita.  Qualquer coisa que não
seja apenas uma letra isolada deve ser da forma ``$[t \circ s]$'' em que
$s$ e $t$ são, eles próprios, bonitos. Reciprocamente, se $t$ e $s$ são bonitos,
então podemos formar um novo termo bonito pondo um $\circ$ entre eles e
cercando-os por um par de colchetes.  Poderíamos usar essas operações
para \emph{definir} o conjunto dos termos bonitos.  Esta é uma \emph{definição
  indutiva}.

\begin{defn}[Termos bonitos]
  O conjunto dos \emph{termos bonitos} é definido indutivamente como segue:
  \begin{enumerate}
  \item Qualquer letra $\mathrm{a}$, $\mathrm{b}$, $\mathrm{c}$,
    $\mathrm{d}$ é um termo bonito.
  \item Se $s_1$ e $s_2$ são termos bonitos, então também o
    é $[s_1 \circ s_2]$.
  \item Nada mais é um termo bonito.
  \end{enumerate}
\end{defn}

Esta definição nos diz que algo conta como um termo bonito sse ele
pode ser construído de acordo com as duas condições (1) e~(2) em algum
número finito de passos. No primeiro passo, construímos todos os termos bonitos
compostos apenas por letras isoladas, isto é,
\[
\mathrm{a}, \mathrm{b}, \mathrm{c}, \mathrm{d}
\]
No segundo passo, aplicamos (2) aos termos que construímos. Obteremos
\[
[\mathrm{a} \circ \mathrm{a}], [\mathrm{a} \circ \mathrm{b}],
[\mathrm{b} \circ \mathrm{a}], \dots, [\mathrm{d} \circ \mathrm{d}]
\]
para todas as combinações de duas letras. No terceiro passo, aplicamos (2)
novamente, a quaisquer dois termos bonitos que construímos até agora. Obtemos novos
termos bonitos como $[\mathrm{a} \circ [\mathrm{a} \circ
    \mathrm{a}]]$---em que $t$ é $\mathrm{a}$ do passo~1 e $s$ é
$[\mathrm{a} \circ \mathrm{a}]$ do passo~2---e $[[\mathrm{b} \circ
    \mathrm{c}] \circ [\mathrm{d} \circ \mathrm{b}]]$ construído a partir
dos dois termos $[\mathrm{b} \circ \mathrm{c}]$ e $[\mathrm{d}
  \circ \mathrm{b}]$ do passo~2. E assim por diante.  A cláusula (3) impede
que qualquer coisa não construída desse modo se infiltre no conjunto dos termos
bonitos.

Note que ainda não provamos que toda sequência de símbolos que
``parece'' bonita é bonita de acordo com esta definição. No entanto,
deve ficar claro que tudo o que podemos construir de fato ``parece
bonito'': os colchetes estão equilibrados, e $\circ$ conecta partes que são,
elas próprias, bonitas.

A característica-chave das definições indutivas é que, se você quer provar
algo sobre todos os termos bonitos, a definição lhe diz quais casos
você deve considerar.  Por exemplo, se lhe é dito que $t$ é um termo
bonito, a definição indutiva lhe diz com o que $t$ pode se parecer: $t$
pode ser uma letra, ou pode ser $[s_1 \circ s_2]$ para algum par de
termos bonitos $s_1$ e~$s_2$. Por causa da cláusula (3), essas são as únicas
possibilidades.

Ao provar afirmações sobre todo um conjunto definido indutivamente, a
forma forte da indução torna-se particularmente importante. Por exemplo,
suponha que queiramos provar que, para todo termo bonito de comprimento~$n$, o
número de $[$ nele é~$< n/2$.  Isso pode ser visto como uma afirmação sobre
todos os~$n$: para todo $n$, o número de $[$ em qualquer termo bonito de
comprimento~$n$ é $< n/2$.

\begin{prop}
  Para qualquer $n$, o número de $[$ em um termo bonito de comprimento~$n$ é
  $< n/2$.
\end{prop}

\begin{proof}
Para provar esse resultado por indução (forte), temos que mostrar que a
seguinte afirmação condicional é verdadeira:
\begin{quote}
  Se, para todo $l < k$, qualquer termo bonito de comprimento~$l$ tem $< l/2$
  $[$, então qualquer termo bonito de comprimento~$k$ tem $< k/2$ $[$.
\end{quote}
Para mostrar este condicional, suponha que seu antecedente é verdadeiro, isto é,
suponha que, para qualquer $l<k$, termos bonitos de comprimento~$l$ contêm $< l/2$
$[$.  Chamamos essa suposição de hipótese indutiva. Queremos
mostrar que o mesmo vale para termos bonitos de comprimento~$k$.

Então suponha que $t$ é um termo bonito de comprimento~$k$.  Como os termos bonitos são
definidos indutivamente, temos dois casos: (1)~$t$ é uma letra
isolada, ou (2)~$t$ é $[s_1 \circ s_2]$ para alguns termos bonitos $s_1$ e~$s_2$.
\begin{enumerate}
\item $t$ é uma letra.  Então $k = 1$, e o número de $[$ em $t$
é~$0$. Como $0 < 1/2$, a afirmação vale.
\item $t$ é $[s_1 \circ s_2]$ para alguns termos bonitos $s_1$ e~$s_2$.
  Sejam $l_1$ o comprimento de~$s_1$ e $l_2$ o comprimento
  de~$s_2$.  Então o comprimento~$k$ de $t$ é $l_1+l_2+3$ (os comprimentos de
  $s_1$ e $s_2$ mais três símbolos $[$, $\circ$, $]$). Como
  $l_1+l_2+3$ é sempre maior que $l_1$, $l_1 < k$. De modo semelhante, $l_2
  < k$. Isso significa que a hipótese de indução se aplica aos termos
  $s_1$ e~$s_2$: o número~$m_1$ de $[$ em $s_1$ é $< l_1/2$, e
  o número~$m_2$ de $[$ em~$s_2$ é $< l_2/2$.

  O número de $[$ em $t$ é o número de $[$ em~$s_1$, mais o
  número de $[$ em~$s_2$, mais~$1$, isto é, é $m_1 + m_2 + 1$. Como $m_1
  < l_1/2$ e $m_2 < l_2/2$, temos:
  \[
  m_1 + m_2 + 1 < \frac{l_1}{2} + \frac{l_2}{2} + 1 = \frac{l_1+l_2+2}{2} < \frac{l_1+l_2+3}{2} = k/2.
  \]
\end{enumerate}
Em cada caso, mostramos que o número de $[$ em $t$ é $< k/2$ (com
base na hipótese indutiva). Por indução forte, a
proposição se segue.
\end{proof}

\begin{prob}
  Defina o conjunto dos termos superbonitos por
  \begin{enumerate}
  \item Qualquer letra $\mathrm{a}$, $\mathrm{b}$, $\mathrm{c}$,
    $\mathrm{d}$ é um termo superbonito.
  \item Se $s$ é um termo superbonito, então também o é $[s]$.
  \item Se $s_1$ e $s_2$ são termos superbonitos, então também o
    é $[s_1 \circ s_2]$.
  \item Nada mais é um termo superbonito.
  \end{enumerate}
  Mostre que o número de $[$ em um termo superbonito~$t$ de comprimento~$n$ é
    $\le n/2 +1$.
\end{prob}

\end{document}
